\section{Background and Related Work}\label{sec:background}

\paragraph{Package dependency networks.}
Empirical studies of npm, PyPI, CRAN, RubyGems and Cargo describe dependency networks that grow
super-linearly. In these networks a small number of packages accumulate very many transitive dependents,
and trivial packages are pervasive~\citep{kikas2017structure,decan2019empirical,abdalkareem2017trivial}.
The 2016 removal of \pkg{left-pad} showed how one tiny package can break a large part of the
ecosystem~\citep{npm2016leftpad}. Dependents constrain the versions they accept with SemVer
ranges~\citep{prestonwerner2013semver}. Caret (\texttt{\^{}}) and tilde (\texttt{\~{}}) ranges dominate
in npm, so new releases propagate automatically~\citep{decan2021semver,pinckney2023semver}. This
permissiveness speeds up the uptake of security fixes~\citep{pinckney2023semver}, and it equally speeds
up the uptake of malicious releases.

\paragraph{Supply-chain attacks on npm.}
\citet{zimmermann2019smallworld} showed that compromising a handful of maintainer accounts would let an
attacker reach a large fraction of npm. Taxonomies of supply-chain attacks~\citep{ohm2020backstabber,
ladisa2023sok} catalogue the main techniques:
\begin{itemize}
  \item typosquatting~\citep{taylor2020typosquatting};
  \item weak maintainer accounts~\citep{zahan2022weaklinks};
  \item abuse of install-time lifecycle scripts~\citep{wyss2022latch}.
\end{itemize}
Notable incidents include \pkg{event-stream}~\citep{npm2018eventstream} and
\pkg{ua-parser-js}~\citep{github2021uaparser}. In September 2025 the \pkg{chalk}/\pkg{debug} account
takeover~\citep{aikido2025chalkdebug} and the s1ngularity/Nx compromise~\citep{wiz2025singularity}
preceded two waves of self-propagating malware.

\paragraph{The Shai-Hulud worms.}
Both waves followed the same pattern. A malicious package version ran a lifecycle script on the machine
that installed it. The script scanned for credentials, exfiltrated them to public GitHub repositories,
and used any npm token it found to republish the token owner's packages with the payload
injected~\citep{stepsecurity2025shaihulud,socket2025shaihulud,unit42025shaihulud,cisa2025npm}.
\begin{itemize}
  \item \textbf{Wave~1} (September 2025) capped republishing at 20 packages per token and compromised
        about \ValOneListed{} packages.
  \item \textbf{Wave~2} (November 2025) used a \pkg{preinstall} hook and a cap of 100, and compromised
        about \ValTwoListed{}~\citep{wiz2025shaihulud2,datadog2025shaihulud2}.
\end{itemize}
Vendor analyses agree that the worm spread from maintainer to maintainer through installs, but
generated most compromised packages by republishing with stolen tokens. As \citet{socket2025shaihulud}
put it, the worm ``cannot self-propagate without those credentials''.

\paragraph{Defences.}
Three families of ecosystem defences are being deployed:
\begin{itemize}
  \item \textbf{Lockfiles and pinning} freeze resolved versions. They delay the uptake of malicious and
        of benign releases alike~\citep{pinckney2023semver,zerouali2022vulnerabilities}.
  \item \textbf{Blocking lifecycle scripts} removes the install-time execution that both waves relied
        on. Options include npm's \texttt{--ignore-scripts}, pnpm's default since version~10, and
        LavaMoat's allow-lists~\citep{npmdocs2026ignorescripts,pnpm2025v10,lavamoat2026allowscripts}.
  \item \textbf{Provenance and trusted publishing} bind releases to CI identities, so a stolen
        long-lived token is no longer enough to publish~\citep{github2023provenance,newman2022sigstore,
        github2025trustedpublishing,slsa2023spec}.
\end{itemize}
After Shai-Hulud, GitHub announced token-lifetime limits, revoked classic tokens and expanded trusted
publishing~\citep{github2025npmplan,github2025classictokens}. Compromised CI pipelines can still publish
with valid provenance~\citep{wiz2025singularity,unit422026npmlandscape}, so these defences are partial.

\paragraph{Epidemics and immunization on networks.}
Epidemic models have been used for computer viruses since \citet{kephart1991directed} placed them on
directed sharing graphs. They were later applied to internet worms such as Code Red~\citep{zou2002codered}.
We build on four strands of theory:
\begin{itemize}
  \item \textbf{SIR dynamics.} The model goes back to \citet{kermack1927contribution}; on networks,
        final outbreak sizes map to bond percolation~\citep{newman2002spread}.
  \item \textbf{Epidemic thresholds.} The threshold vanishes on scale-free networks with diverging
        degree moments~\citep{pastorsatorras2001epidemic,pastorsatorras2015review}. On general graphs it
        is governed by the spectral radius $\lambda_{\max}$ of the contagion
        matrix~\citep{wang2003eigenvalue,chakrabarti2008thresholds}. On directed networks $1/\lambda_{\max}$
        is only a lower bound, and directionality raises the true
        threshold~\citep{li2013directed,boguna2005directed}.
  \item \textbf{Immunization.} Targeting high-degree nodes is far more effective than random
        immunization~\citep{pastorsatorras2002immunization}. Acquaintance immunization approximates
        targeting without global knowledge~\citep{cohen2003acquaintance}. Betweenness-based removal
        fragments networks quickly~\citep{holme2002attack,freeman1977betweenness}, although optimized
        strategies beat highest-betweenness heuristics~\citep{schneider2011limited,chen2008immunization}.
  \item \textbf{Methods.} We fit heavy tails with the method of~\citet{clauset2009powerlaw}, as
        implemented by~\citet{alstott2014powerlaw}. Strictly scale-free real networks are
        rare~\citep{broido2019rare}.
\end{itemize}

\paragraph{Positioning.}
Prior work on propagation in package ecosystems is mostly about the \emph{static reachability} of known
vulnerabilities~\citep{decan2018vulnerabilities,liu2022demystifying,zerouali2022vulnerabilities,
mir2023maven,robinson2026originalsin}. The closest dynamic model, by \citet{schueller2024risks}, couples
dependency and shared-maintainer failures in Rust, but it models non-adversarial maintainer exit. To our
knowledge, no prior work models self-propagating npm worms as an epidemic process with a credential
channel, calibrates such a model to observed incidents, or evaluates ecosystem defences against it.
